\subsection{The Direct (Forward‐Mode) Method}  \label{sec:direct}


Differentiating~\eqref{eq:equil} in the direction
$\dot{\HostParameters{}}\in\mathbb{R}^{n_p}$ yields the \emph{tangent system}
%
\begin{equation}\label{eq:forward}
  \bigl(\partial_{\bm{u}}\bm{g}\bigr)\,\dot{\bm{u}}
  \;=\;-\bigl(\partial_{\HostParameters{}}\bm{g}\bigr)\,\dot{\HostParameters{}},
  \qquad
  \dot{\bm{u}}\equiv\dfrac{\partial \bm{u}}{\partial \HostParameters{}}\dot{\HostParameters{}}.
\end{equation}
%
The linear map
$\dot{\HostParameters{}}\mapsto\dot{\bm{u}}$ is the \emph{Jacobian–vector product} (JVP)
%
\[
 \dot{\TangentState} = \mathrm{JVP}_{\bm{h}}(\HostParameters{};\dot{\HostParameters{}})
  \;=\;
  \partial\bm{h}(\HostParameters{})\,\dot{\HostParameters{}},
\]
%
and its numerical evaluation is called forward-mode automatic
differentiation in the \autodiff{} literature\footnote{The \autodiff{} literature calls this the \emph{discrete–direct} or
\emph{tangent‐linear} approach \citep{barthelemy1995automatic}.}.  
Substituting $\dot{\bm{u}}$
into:
%
\[
  d\Objective
  \;=\;
  \bigl\langle\partial_{\bm{u}} \Objective,\dot{\bm{u}}\bigr\rangle
  +\bigl\langle\partial_{\HostParameters{}} \Objective,\dot{\HostParameters{}}\bigr\rangle
\]
%
provides the directional derivative $d \Objective(\HostParameters{};\dot{\HostParameters{}})$ where 
$\partial_{\bm{u}} \Objective$ and $\partial_{\HostParameters{}} \Objective$ are column gradients of a scalar
functional $\Objective :\mathbb{R}^{n_u}\!\times\!\mathbb{R}^{n_p}\to\mathbb{R}$.

\paragraph{Summary.}
Given a nonlinear residual equation $\bm{g}(\bm{u},\HostParameters{})=\bm{0}$, the \emph{direct differentiation
method} (DDM) forms its tangent (Jacobian) system and \emph{propagates} a
perturbation of the parameters
$\dot{\HostParameters{}}\in\mathbb{R}^{n_p}$ forward to the state perturbation
$\dot{\bm{u}}\in\mathbb{R}^{n_u}$ by solving \eqref{eq:forward}.
%
% \begin{equation}\label{eq:forward}
% \frac{\partial\bm{g}}{\partial\bm{u}}\;\dot{\bm{u}}
%  \;=\; -\frac{\partial\bm{g}}{\partial\HostParameters{}}\;\dot{\HostParameters{}} .
% \end{equation}
%
Subsequently, gradients of any scalar objective
$\ell(\bm{u},\HostParameters{})$ follow from the total derivative
$d \ell = \ell_{\bm{u}}\dot{\bm{u}} + \ell_{\HostParameters{}}\dot{\HostParameters{}}$.



\paragraph{Complexity.}
Solving~\eqref{eq:forward} once costs roughly as much as one Newton step of the
primal problem and yields one column of $\partial\bm{h}$.  Hence the work scales
like $\mathcal{O}(n_p)$, whereas memory is $\,\mathcal{O}(1)$ with respect to the
depth of the nonlinear solve.

The cost scales with the \emph{number of inputs}
$n_p$ because each independent direction
$\dot{\HostParameters{}}$ requires a solve of \eqref{eq:forward}.
Memory is modest: no intermediate states need to be stored.